% Erdős Problem #699 (Erdős-Szekeres binomial-gcd conjecture) -- partial resolution.
% Author (per document): "GPT 5.6 Sol Pro", prompted by Liam Price.
% Source: linked from https://www.erdosproblems.com/699 (page edited 2026-07-19);
%   Overleaf: https://www.overleaf.com/read/fptssppkmgpr
% Retrieved and filed 2026-07-24. EXTERNAL work (not ours), kept for reference.
% Proves the conjecture for j <= 3i/2 and for the central case n = 2j.
% NB: subsumes our own "i >= n/5" fragment -- its Lemma 2.7 uses the identical
%     prime-in-(n-i,n] mechanism (q > n/2, n-q < i => v_q(C(n,r))=1 on [i,n/2]).
% Verbatim copy; do not edit the mathematics.

\documentclass[a4paper,11pt]{amsart}

\usepackage[margin=1in]{geometry}
\usepackage[T1]{fontenc}
\usepackage{amsmath,amssymb}
\usepackage{booktabs}
\usepackage{microtype}

\newtheorem{theorem}{Theorem}[section]
\newtheorem{proposition}[theorem]{Proposition}
\newtheorem{lemma}[theorem]{Lemma}
\newtheorem{corollary}[theorem]{Corollary}
\theoremstyle{remark}
\newtheorem{remark}[theorem]{Remark}

\newcommand{\Exc}{\mathcal E}
\newcommand{\lcmop}{\operatorname{lcm}}

\title{Common Prime Divisors of Binomial Coefficients}
\author{GPT 5.6 Sol Pro}
\date{}
\subjclass[2020]{Primary 11B65; Secondary 11A41}
\keywords{binomial coefficient, common prime divisor, Kummer's theorem}

\begin{document}

\begin{abstract}
Erd\H{o}s and Szekeres conjectured that, whenever
$1\le i<j\le n/2$, the binomial coefficients $\binom{n}{i}$ and
$\binom{n}{j}$ have a common prime divisor at least $i$.
We prove the conjecture when $j\le 3i/2$ and in the central case
$n=2j$.
The proof localizes the prime-power divisors of $\binom{n}{i}$ supported
on primes at least $i$ and combines the resulting divisibility with a
theorem of Ecklund, Eggleton, Erd\H{o}s, and Selfridge.
For odd prime indices in the central case, the argument uses an analytic
estimate for the infinite range and an exact finite computation whose
complete code is included.
\end{abstract}

\maketitle

\section{Introduction}

Erd\H{o}s and Szekeres \cite{ErSz} conjectured that, for integers
$1\le i<j\le n/2$, there is a prime $q\ge i$ such that
$q\mid\gcd\bigl(\binom{n}{i},\binom{n}{j}\bigr)$. We establish the
following two ranges.

\begin{theorem}\label{thm:main}
Let $1\le i<j\le n/2$. If either $j\le 3i/2$ or $n=2j$, then some
prime $q\ge i$ divides both $\binom{n}{i}$ and $\binom{n}{j}$.
\end{theorem}

All logarithms are natural. For a prime $q$, let $v_q$ denote the
$q$-adic valuation, and let $\pi(x)$ denote the number of primes not
exceeding $x$. Define the part of $\binom{n}{i}$ supported on primes at
least $i$ by
\[
 V_i(n)=\prod_{\substack{q\ge i\\q\text{ prime}}}
 q^{v_q(\binom{n}{i})},
 \qquad U_i(n)=\binom{n}{i}/V_i(n).
\]
Thus every prime divisor of $U_i(n)$ is less than $i$, whereas every
prime divisor of $V_i(n)$ is at least $i$. The only substantial
external result used below is the following theorem, stated here in the
form needed.

\begin{theorem}[Ecklund--Eggleton--Erd\H{o}s--Selfridge]
\label{thm:eees}
Let $N\ge 2r$ be positive integers, and write $\binom{N}{r}=uv$, where
all prime divisors of $u$ are less than $r$ and all prime divisors of
$v$ are at least $r$. Then $u<v$, except when
\begin{align*}
\Exc=\{&(8,3),(9,4),(10,5),(12,5),(21,7),(21,8),\\
       &(30,7),(33,13),(33,14),(36,13),(36,17),(56,13)\},
\end{align*}
where each ordered pair is $(N,r)$.
\end{theorem}

This is the main theorem of \cite{EEES}. Its exceptional pairs will not
cause exceptions to Theorem~\ref{thm:main}.

\begin{lemma}\label{lem:exceptions}
If $(n,i)\in\Exc$ and $i<j\le n/2$, then some prime $q\ge i$ divides
both $\binom{n}{i}$ and $\binom{n}{j}$.
\end{lemma}

\begin{proof}
The pairs $(9,4)$ and $(10,5)$ admit no such $j$. For every remaining
pair, a suitable prime is given in the table.
\begin{center}
\begin{tabular}{cc@{\qquad}cc}
\toprule
$(n,i)$ & $q$ & $(n,i)$ & $q$\\
\midrule
$(8,3)$   & $7$  & $(12,5)$  & $11$\\
$(21,7)$  & $19$ & $(21,8)$  & $19$\\
$(30,7)$  & $29$ & $(33,13)$ & $31$\\
$(33,14)$ & $31$ & $(36,13)$ & $31$\\
$(36,17)$ & $31$ & $(56,13)$ & $53$\\
\bottomrule
\end{tabular}
\end{center}
In each case, $q>n/2$ and $n-q<i$. Hence $n-q<r<q$ for every
$i\le r\le n/2$. Since $q^2>n$, it follows that
$v_q\bigl(\binom{n}{r}\bigr)=1$; apply this with $r=i$ and $r=j$.
\end{proof}

\section{Prime-power localization}

The next lemma is the basic divisibility mechanism.

\begin{lemma}[Transfer and localization]\label{lem:localization}
Let $1\le i<j\le n/2$, put $c=n-j$, and suppose that
$a=v_q\bigl(\binom{n}{i}\bigr)>0$ while
$q\nmid\binom{n}{j}$.

\smallskip
\noindent
\emph{(a)} One has $q^a\mid\binom{j}{i}$ and
$q^a\mid\binom{c}{i}$.

\smallskip
\noindent
\emph{(b)} If $q\ge i$, then there are nonnegative integers $r,s$ with
$r+s<i$ such that $q^a\mid j-r$ and $q^a\mid c-s$.
\end{lemma}

\begin{proof}
The two identities
\[
 \binom{n}{i}\binom{n-i}{j-i}
   =\binom{n}{j}\binom{j}{i},
 \qquad
 \binom{n}{i}\binom{n-i}{c-i}
   =\binom{n}{c}\binom{c}{i}
\]
and $\binom{n}{c}=\binom{n}{j}$ prove part~(a).

Assume $q\ge i$. Among the $i$ integers
$n-i+1,\ldots,n$ there is at most one multiple of $q$. Moreover,
$v_q(i!)=0$ if $q>i$, and $v_q(i!)=1$ if $q=i$. Consequently there is
a unique $t\in\{0,\ldots,i-1\}$ such that $q^a\mid n-t$.

Let $r$ and $s$ be the least nonnegative residues of $j$ and $c$ modulo
$q^a$. If $s_q(N)$ is the sum of the base-$q$ digits of $N$, then
Legendre's formula gives
\[
 v_q\binom{n}{j}
 =\frac{s_q(j)+s_q(c)-s_q(n)}{q-1}.
\]
The right-hand side counts the carries in the base-$q$ addition
$j+c=n$, which is Kummer's theorem in this case. The valuation is zero,
so the first $a$ digits add without a carry out and $r+s<q^a$. Since
$r+s\equiv n\equiv t\pmod{q^a}$ and $0\le t<i\le q\le q^a$, one has
$r+s=t<i$. The definitions of $r$ and $s$ give the required
divisibilities.
\end{proof}

\begin{corollary}\label{cor:triangle}
If no prime $q\ge i$ divides both $\binom{n}{i}$ and $\binom{n}{j}$,
then
\[
 V_i(n)\mid
 \lcmop_{\substack{r,s\ge0\\r+s<i}}
 \gcd(j-r,n-j-s).
\]
In particular, $V_i(n)\mid\binom{j}{i}$ and
$V_i(n)\mid\binom{n-j}{i}$.
\end{corollary}

\begin{proof}
Apply Lemma~\ref{lem:localization} to the full $q$-part of $V_i(n)$ for
each prime $q\ge i$. Part~(b) places that full $q$-part in one of the
greatest common divisors occurring in the least common multiple, and
part~(a) gives the final two divisibilities.
\end{proof}

\section{Proof of Theorem~\ref{thm:main}}

\subsection{The range $j\le3i/2$}

\begin{proposition}\label{prop:near}
The conclusion of Theorem~\ref{thm:main} holds when $j\le3i/2$.
\end{proposition}

\begin{proof}
Suppose otherwise and put $d=j-i$. Corollary~\ref{cor:triangle} gives
$V_i(n)\mid\binom{j}{i}=\binom{j}{d}$. By
Lemma~\ref{lem:exceptions}, $(n,i)\notin\Exc$, so
Theorem~\ref{thm:eees} applies. The inequality $j\le3i/2$ is equivalent
to $3d\le j$, and hence $2d\le i<j$. Vandermonde's convolution and
monotonicity along the first half of row $2j$ now give the contradictory
chain
\[
 \binom{n}{i}<V_i(n)^2\le\binom{j}{d}^{\!2}
 <\binom{2j}{2d}\le\binom{2j}{i}\le\binom{n}{i}.
\]
The strict middle inequality holds because the Vandermonde sum for
$\binom{2j}{2d}$ contains the positive term $\binom{j}{d}^2$ and at
least one other positive term.
\end{proof}

\subsection{The central case}

At the center of the row, the two residues furnished by
Lemma~\ref{lem:localization} must coincide.

\begin{lemma}[Central compression]\label{lem:central}
Let $1\le i<m$, and suppose that no prime $q\ge i$ divides both
$\binom{2m}{i}$ and $\binom{2m}{m}$. If $h=\lceil i/2\rceil$, then
$V_i(2m)\mid\binom{m}{h}$.
\end{lemma}

\begin{proof}
Let $q^a$ be the full $q$-part of $V_i(2m)$. Applying
Lemma~\ref{lem:localization} with $j=c=m$ gives nonnegative $r,s$ such
that $r+s<i$ and $q^a\mid m-r,m-s$. Thus $q\mid r-s$. Since
$|r-s|<i\le q$, one has $r=s$, so $2r<i$ and $0\le r<h$. Hence every
full prime-power part of $V_i(2m)$ divides one of
$m,m-1,\ldots,m-h+1$. Every prime divisor of $V_i(2m)$ is greater than
$h$; therefore none is removed on division by $h!$, and
$V_i(2m)\mid\binom{m}{h}$.
\end{proof}

\begin{proposition}\label{prop:central}
For every $m\ge2$ and $1\le i<m$, some prime $q\ge i$ divides both
$\binom{2m}{i}$ and $\binom{2m}{m}$.
\end{proposition}

\begin{proof}
For $i=1$, the prime $2$ works because
$\binom{2m}{m}=2\binom{2m-1}{m-1}$. Suppose first that $i=2h$ is
even. No pair in $\Exc$ has both coordinates even. Under the contrary
assumption, Theorem~\ref{thm:eees}, Lemma~\ref{lem:central}, and
Vandermonde's convolution give
\[
 \binom{2m}{2h}<V_{2h}(2m)^2\le\binom{m}{h}^{\!2}
 <\binom{2m}{2h},
\]
a contradiction.

Now let $i$ be odd and composite. The even case applied to $i-1$
produces a prime $q\ge i-1$ dividing both $\binom{2m}{i-1}$ and
$\binom{2m}{m}$. Since $i-1>2$ and $i$ are composite, $q>i$. The
identity
$\binom{2m}{i}=\binom{2m}{i-1}(2m-i+1)/i$ then shows that $q$ also
divides $\binom{2m}{i}$.

It remains to take $i=p=2k+1$ prime. The exceptional instances of
Theorem~\ref{thm:eees} are covered by Lemma~\ref{lem:exceptions}, so
assume $(2m,p)\notin\Exc$. If the proposition were false, then
Lemma~\ref{lem:central} and Theorem~\ref{thm:eees} would imply
\begin{equation}\label{eq:Rless}
 R_k(m):=\frac{\binom{2m}{2k+1}}{\binom{m}{k+1}^{\!2}}<1.
\end{equation}
We derive the opposite inequality.

For $0\le r<k$, set $x_r=2m-(2r+1)$ and
$X=\prod_{r=0}^{k-1}x_r$. No prime $q>p$ divides $X$. Indeed, if
$q\mid x_r$, then $q\mid\binom{2m}{p}$ because $q>p$. Let $\rho$ be
the least nonnegative residue of $m$ modulo $q$. Since
$2\rho\equiv2r+1\pmod q$ and $0<2r+1<q$, either
$2\rho=2r+1$ or $2\rho=q+2r+1$; the first equality is impossible by
parity. Thus adding $m+m$ in base $q$ produces a carry in the units
place, and Kummer's theorem gives $q\mid\binom{2m}{m}$, a contradiction.
The same carry argument applies to $q=p$. If $v_p(X)\ge2$, then the
numerator of $\binom{2m}{p}$ contains at least two factors of $p$, while
$p!$ contains exactly one; hence $p$ would divide both binomial
coefficients. Therefore
\begin{equation}\label{eq:support}
 q\nmid X\quad(q>p),\qquad v_p(X)\le1.
\end{equation}

Put $a=2m-2k+1$ and $e=k-\pi(p)+2$. For each odd prime $q<p$, choose
an index $r_q$ at which $v_q(x_r)$ is maximal, and delete every chosen
index. At most $\pi(p)-2$ distinct indices are deleted, so at least
$e$ remain. If $r$ remains, then
$q^{v_q(x_r)}\mid x_r-x_{r_q}=2(r_q-r)$ and, as $q$ is odd,
$q^{v_q(x_r)}\mid r-r_q$. Thus, for the set $S$ of remaining indices,
\[
 q^{\sum_{r\in S}v_q(x_r)}
 \mid\prod_{r\in S}|r-r_q|
 \mid r_q!(k-1-r_q)!\mid(k-1)!.
\]
Together with \eqref{eq:support}, this shows that
$\prod_{r\in S}x_r\mid p(k-1)!$. Since $x_r\ge a$ and $|S|\ge e$,
we obtain
\begin{equation}\label{eq:smooth}
 a^e\le p(k-1)!.
\end{equation}

A direct calculation gives
\begin{equation}\label{eq:Rmonotone}
 \frac{R_k(m+1)}{R_k(m)}
 =\frac{(2m+1)(m-k)}{(2m-2k+1)(m+1)}<1,
\end{equation}
where the denominator exceeds the numerator by $2m-k+1$. Thus
$R_k(m)$ decreases with $m$.

We first treat $k\ge2500$. Put $K=2500$, $P=5001$, and
\[
 C=\left(2+\frac1K\right)
   \left(1+\frac{3}{2\log P}\right).
\]
The explicit estimate
\[
 \pi(x)<\left(1+\frac{3}{2\log x}\right)\frac{x}{\log x}
 \qquad(x>1),
\]
recorded in \cite[(2)]{EEES}, gives $\pi(p)/k<C/\log p$. Using
$(k-1)!\le k^{k-1}$ in \eqref{eq:smooth}, the fact that
$x\mapsto\log(2x+1)/x$ decreases, and $\log k<\log p$, we obtain
\[
 \log\frac ak
 \le\frac{\log p+(\pi(p)-3)\log k}{k-\pi(p)+2}
 <\frac{C+\log P/K}{1-C/\log P}.
\]
The remaining numerical comparison can be made with rational bounds.
The exponential series gives $8/3<\exp(1)<49/18<11/4$, and
$(11/4)^8<5001<(8/3)^9$; hence $8<\log P<9$. Consequently
\[
 \frac{C+\log P/K}{1-C/\log P}
 <\frac{95163/40000}{224981/320000}
 =\frac{761304}{224981}<\frac{17}{5}<\log32.
\]
For the last inequality, the power series
$\log2=2\sum_{\nu\ge0}((2\nu+1)3^{2\nu+1})^{-1}$ gives
$\log32=5\log2>280/81>17/5$. It follows that $a<32k$, and hence
$m<17k$.

Vandermonde's convolution gives
\[
 R_k(m)=2\sum_{s=0}^{k}
 \frac{\binom{m}{k-s}\binom{m}{k+1+s}}
      {\binom{m}{k+1}^{\!2}},
\]
and the summand without the leading factor $2$ equals
\[
 \frac{k+1}{m-k}
 \prod_{u=1}^{s}
 \frac{(k+1-u)(m-k-u)}{(k+1+u)(m-k+u)}.
\]
Let $L=\lfloor\sqrt{k}/4\rfloor$. For $0\le s\le L$, the first factor
is greater than $1/16$, while every factor in the product is at least
$(1-2L/k)^2$. Bernoulli's inequality and $L^2\le k/16$ give
\[
 \frac{\binom{m}{k-s}\binom{m}{k+1+s}}
      {\binom{m}{k+1}^{\!2}}
 >\frac1{16}\left(1-\frac{2L}{k}\right)^{2L}
 \ge\frac1{16}\left(1-\frac{4L^2}{k}\right)
 \ge\frac3{64}.
\]
Since $L\ge12$, this yields
$R_k(m)>2(L+1)(3/64)\ge39/32>1$.

It remains to consider $k<2500$. Let $A_k$ be the largest odd integer
such that $A_k^e\le p(k-1)!$, and put
$M_k=(A_k+2k-1)/2$. Inequality \eqref{eq:smooth} gives $m\le M_k$,
while \eqref{eq:Rmonotone} reduces the required inequality to
$R_k(M_k)>1$ whenever $M_k\ge p+1$. Appendix~\ref{app:code} verifies
all these inequalities by exact integer arithmetic. Thus
$R_k(m)>1$ in every case, contradicting \eqref{eq:Rless} and completing
the proof.
\end{proof}

Propositions~\ref{prop:near} and \ref{prop:central} prove
Theorem~\ref{thm:main}.

\section{A Euclidean-division bound}

The localization argument also gives a necessary condition for any
counterexample outside the two ranges of Theorem~\ref{thm:main}. For
$T\ge1$, define
\[
 \Lambda_i(T)=\prod_{\substack{q\ge i\\q\text{ prime}}}
 q^{\lfloor\log_qT\rfloor},
 \qquad \Lambda_i(0)=1.
\]
Thus $\Lambda_i(T)$ is the part of $\lcmop(1,\ldots,T)$ supported on
primes at least $i$.

\begin{proposition}\label{prop:euclidean}
Suppose $1\le i<j\le n/2$ and no prime $q\ge i$ divides both
$\binom{n}{i}$ and $\binom{n}{j}$. Write $n=Qj+R$, where
$Q=\lfloor n/j\rfloor\ge2$ and $0\le R<j$, and put
\[
 T=R+(Q-1)(i-1),
 \qquad
 h=\max\left\{0,\left\lceil\frac{i-R}{Q}\right\rceil\right\}.
\]
Then $V_i(n)\mid\Lambda_i(T)\binom{j}{h}$, and consequently
\[
 \binom{n}{i}<\left(\Lambda_i(T)\binom{j}{h}\right)^2.
\]
\end{proposition}

\begin{proof}
Put $c=n-j=(Q-1)j+R$. For each full $q$-part $q^a$ of $V_i(n)$,
Lemma~\ref{lem:localization} supplies $r,s\ge0$ such that $r+s<i$ and
$q^a\mid j-r,c-s$. Hence
$q^a\mid E:=R+(Q-1)r-s$. If $E\ne0$, then
\[
 |E|\le\max\{R+(Q-1)(i-1),i-1-R\}\le T,
\]
so $q^a\mid\Lambda_i(T)$. If $E=0$, then
$s=R+(Q-1)r$ and $R+Qr<i$, whence $0\le r<h$ and $q^a\mid j-r$.
The product of all full prime-power parts arising in this second way
divides $j(j-1)\cdots(j-h+1)$. Every prime involved is greater than
$h$---for $i\ge2$, one has $h\le\lceil i/2\rceil<i$, while for $i=1$
one has $h\le1<q$---so this product divides $\binom{j}{h}$. Multiplying
the two classes of prime-power parts proves the first assertion.

A counterexample cannot have $(n,i)\in\Exc$ by
Lemma~\ref{lem:exceptions}. Theorem~\ref{thm:eees} therefore gives
$\binom{n}{i}<V_i(n)^2$, and the stated inequality follows.
\end{proof}

\appendix
\section{Exact finite verification}\label{app:code}

The following Python 3 program uses only exact integer arithmetic. It
enumerates every $k<2500$ for which $p=2k+1$ is prime, computes the
endpoint $M_k$ used in the proof of Proposition~\ref{prop:central}, and
asserts $R_k(M_k)>1$ whenever $p<m\le M_k$ is nonempty. There are
$666$ such endpoint checks. The smallest verified endpoint ratio is
attained at $(k,p,e,A_k,M_k)=(9,19,3,91,54)$ and is
\[
 \frac{659645433004618405800}{572679032824811448900}>1.
\]
The program verifies these data as a reproducibility certificate.

\begin{verbatim}
from math import comb, isqrt

K0 = 2500
limit = 2 * K0 + 1

prime = [True] * (limit + 1)
prime[0] = prime[1] = False
for q in range(2, isqrt(limit) + 1):
    if prime[q]:
        for z in range(q * q, limit + 1, q):
            prime[z] = False

pi = [0] * (limit + 1)
count = 0
for x in range(limit + 1):
    if prime[x]:
        count += 1
    pi[x] = count

def largest_odd_root(N, e):
    """Largest odd A such that A**e <= N."""
    lo, hi = 0, 1
    while hi**e <= N:
        hi *= 2
    while lo + 1 < hi:
        mid = (lo + hi) // 2
        if mid**e <= N:
            lo = mid
        else:
            hi = mid
    return lo if lo % 2 else lo - 1

checked = 0
factorial = 1                 # (k - 1)! when k = 1
best = None

for k in range(1, K0):
    if k > 1:
        factorial *= k - 1

    p = 2 * k + 1
    if not prime[p]:
        continue

    e = k - pi[p] + 2
    A = largest_odd_root(p * factorial, e)
    M = (A + 2 * k - 1) // 2

    if M < p + 1:
        continue

    lhs = comb(2 * M, p)
    rhs = comb(M, k + 1)**2
    assert lhs > rhs

    datum = (lhs, rhs, k, p, e, A, M)
    if best is None or lhs * best[1] < best[0] * rhs:
        best = datum
    checked += 1

certificate = (
    659645433004618405800,
    572679032824811448900,
    9, 19, 3, 91, 54
)

assert checked == 666
assert best == certificate
print(checked)
print(best)
\end{verbatim}

The output is
\begin{verbatim}
666
(659645433004618405800, 572679032824811448900,
 9, 19, 3, 91, 54)
\end{verbatim}

\begin{thebibliography}{9}

\bibitem{EEES}
E.~F. Ecklund, Jr., R.~B. Eggleton, P.~Erd\H{o}s, and J.~L. Selfridge,
\emph{On the prime factorization of binomial coefficients},
J. Austral. Math. Soc. Ser. A \textbf{26} (1978), 257--269.

\bibitem{ErSz}
P. Erd\H{o}s and G. Szekeres,
\emph{Some number theoretic problems on binomial coefficients},
Austral. Math. Soc. Gaz. \textbf{5} (1978), 97--99.

\end{thebibliography}

\end{document}